% ============================================================
% preambule.tex -- commun a tout le memoire.
% Compile avec pdflatex (Overleaf : babel francais) ou xelatex/lualatex (polyglossia).
% ============================================================
\usepackage{iftex}
\ifPDFTeX
  \usepackage[utf8]{inputenc}
  \usepackage[T1]{fontenc}
  \usepackage{mathpazo}
  \usepackage[french]{babel}
\else
  \usepackage{fontspec}
  \usepackage{polyglossia}
  \setmainlanguage{french}
  \setmainfont{TeX Gyre Pagella}
  \setmonofont{DejaVu Sans Mono}[Scale=0.85]
\fi
\usepackage{amsmath,amssymb,amsthm}
\usepackage{microtype}
\usepackage{graphicx}
\usepackage{booktabs,longtable,array,tabularx}
\usepackage{caption}
\usepackage{enumitem}
\usepackage{placeins}
\usepackage[table]{xcolor}
\usepackage{tikz}
\usetikzlibrary{arrows.meta,positioning,calc,patterns,decorations.pathreplacing}
\usepackage{pgfplots}\pgfplotsset{compat=1.16}
\usepackage[a4paper,margin=2.6cm]{geometry}
\usepackage{fancyhdr}
\usepackage{titlesec}
\emergencystretch=3em
\setlength{\parskip}{0.4em}
\setlength{\parindent}{0pt}

% ---- palette
\definecolor{navy}{HTML}{0B1D3A}
\definecolor{steel}{HTML}{3A6EA5}
\definecolor{ice}{HTML}{E7EEF6}
\definecolor{alert}{HTML}{D62839}
\definecolor{ok}{HTML}{2E8B57}
\definecolor{muted}{HTML}{6B7B8C}
\definecolor{amber}{HTML}{B7791F}

\usepackage[colorlinks=true,linkcolor=steel,citecolor=steel,urlcolor=steel]{hyperref}

\setcounter{secnumdepth}{3}\setcounter{tocdepth}{3}
\titleformat{\chapter}[display]{\normalfont\huge\bfseries\color{navy}}{\chaptertitlename\ \thechapter}{14pt}{\Huge}
\titleformat{name=\chapter,numberless}[block]{\normalfont\Large\scshape\color{navy}\raggedleft}{}{0pt}{}
\titlespacing*{\chapter}{0pt}{60pt}{36pt}
\titleformat{\section}{\normalfont\Large\bfseries\color{steel}}{\thesection}{0.8em}{}
\titleformat{\subsection}{\normalfont\large\bfseries\color{navy}}{\thesubsection}{0.8em}{}
\titleformat{\subsubsection}{\normalfont\normalsize\bfseries\color{navy}}{\thesubsubsection}{0.8em}{}
\newcommand{\usection}[1]{\section*{#1}\addcontentsline{toc}{section}{#1}}
\pagestyle{fancy}\fancyhf{}
\renewcommand{\chaptermark}[1]{\markboth{\textsc{#1}}{}}
\renewcommand{\sectionmark}[1]{\markright{#1}}
\fancyhead[L]{\small\color{navy}\nouppercase{\leftmark}}
\fancyhead[R]{\small\color{navy}\nouppercase{\rightmark}}
\fancyfoot[C]{\thepage}
\renewcommand{\headrulewidth}{0.5pt}
\renewcommand{\headrule}{{\color{navy}\hrule width\headwidth height\headrulewidth}}
\fancypagestyle{plain}{\fancyhf{}\fancyfoot[C]{\thepage}\renewcommand{\headrulewidth}{0pt}}

% ---- code et statuts epistemiques
\newcommand{\code}[1]{\texttt{\detokenize{#1}}}
\newcommand{\F}{\textcolor{ok}{\textbf{[F]}}}      % fait lu dans le code ou les mesures deposees
\newcommand{\I}{\textcolor{amber}{\textbf{[I]}}}   % inference
\newcommand{\Hy}{\textcolor{alert}{\textbf{[H]}}}  % hypothese non testee
\newcommand{\Pc}{\textcolor{ok}{\textbf{[C]}}}     % propriete justifiee par le code
\newcommand{\Pe}{\textcolor{amber}{\textbf{[E]}}}  % propriete verifiee experimentalement
\newcommand{\ie}{c.-à-d.\ }
\newcommand{\gf}{\mathbb{F}_{16}}
\newcommand{\vpshufb}{\code{vpshufb}}

% ---- theoremes
\theoremstyle{plain}
\newtheorem{proposition}{Proposition}[chapter]
\newtheorem{lemme}[proposition]{Lemme}
\newtheorem{corollaire}[proposition]{Corollaire}
\theoremstyle{definition}
\newtheorem{definition}[proposition]{Définition}
\theoremstyle{remark}
\newtheorem{remarque}[proposition]{Remarque}

\usepackage{framed}
\colorlet{shadecolor}{ice}
\newenvironment{encadre}{\begin{shaded}\small}{\end{shaded}}

\usepackage{seqsplit}
\newcommand{\sha}[1]{{\ttfamily\scriptsize\seqsplit{#1}}}
